\chapter{Conclusioni}
%\addcontentsline{toc}{section}{\textbf{Conclusioni}}

Il 29 luglio 2025 un terremoto di magnitudo \(M8.8\) ha colpito le coste della Kamchatka nella Russia orientale. A seguito dell'evento è stata registrata una ripresa di attività eruttiva nei vulcani Klyuchevskoy e Krasheninnikov, situati a decine di chilometri di distanza dall'epicentro del sisma.
Utilizzando un modello a faglia finita basato su dislocazioni rettangolari, sono stati calcolati gli spostamenti superficiali indotti dal terremoto. La forma del campo di deformazione superficiale è stata confrontata con quella ottenuta dall'ente statunitense U.S. Geological Survey, con cui risulta in accordo. 
Con lo stesso modello si è calcolato il campo di sforzo generato dal terremoto, per verificare se questo fosse in grado di modificare il campo regionale già presente. Per riprodurre fedelmente il campo regionale è stato adottato un modello semplificato, basato su considerazioni tettoniche e sugli effetti della topografia, implementata con una geometria semplificata.
I risultati mostrano che in prossimità dei vulcani, tra i \qty{200}{\kilo\meter} e i \qty{400}{\kilo\meter} dall'epicentro, il campo di sforzo regionale (\qty{1e6}{\pascal} - \qty{1e7}{\pascal}) è decisamente più intenso di quello del terremoto (\qty{1e4}{\pascal}), di almeno due o tre ordini di grandezza. Conseguentemente, l'aggiunta del campo del terremoto non è sufficiente da indurre variazioni nell'orientazione degli assi principali degli sforzi regionali.
Questo fatto si rilette nelle traiettorie di dicchi magmatici, che si orientano perpendicolarmente allo sforzo meno compressivo. Simulando le traiettorie dei dicchi si è visto che questi non subiscono variazioni significative delle loro traiettorie, spostandosi al più di \qty{1e-5}{\kilo\meter}.
Si può stabilire che la variazione del campo di sforzo statico non sia stata la causa che ha portato all'eruzione dei due vulcani. Nonostante ciò, il modello utilizzato può essere testato nella zona di massima intensità degli spostamenti superficiali, per verificare eventuali cambiamenti nel campo di sforzo che potrebbero favorire vulcanismo dai due ai cinque anni successivi all'evento.
Non va rigettata nemmeno l'ipotesi che sia stato il terremoto a indurre le eruzioni. In questa tesi si è studiato un modello basato sul campo di sforzo statico e sulle sue deformazioni. Non è stato verificato un approccio dinamico, basato sulle oscillazioni indotte dalle onde sismiche, che potrebbero aver innescato un processo, già predisposto, di eruzione.